\documentclass{article}
\usepackage[UTF8]{ctex}
\usepackage{amsmath,amssymb,amsfonts}
\usepackage{enumitem}
\usepackage{geometry}
\geometry{a4paper, margin=1in}
% === 错题本排版优化（不参与题目标识） ===
\setlength{\parindent}{0pt}
\setlength{\parskip}{0.35em plus 0.10em minus 0.08em}
\setlength{\emergencystretch}{3em}
\setlist{itemsep=0.25em, topsep=0.40em, parsep=0pt}
\allowdisplaybreaks[2]
% === 错题本排版优化结束 ===
\begin{document}
\section*{浙江大学2026年高等代数真题}
\section*{96．浙江大学2026年高等代数试卷}
一．填空题。\\

\subsection{题目1}
1．假如 \(A\) 是 4 阶整数矩阵，其有特征值 \(\sqrt{5}+\sqrt{3}\) ，那么在复数范围内 \(A\) 的相似标准型为 \(\_\_\_\_\)，\(\left|2 \sqrt{5} E-A^{*}\right|=\) \(\_\_\_\_\)。

\subsection{题目2}
2．设二次型 \(f\left(x_{1}, x_{2}, x_{3}\right)=2 x_{1}^{2}+2 t x_{1} x_{2}+4 x_{1} x_{3}+x_{2}^{2}+3 x_{3}^{2}\) ，当 \(t\) 满足 \(\_\_\_\_\)时，\(f\) 是正定的，当\(t\) 满足 \(\_\_\_\_\)时，\(f\) 的负惯性指数是 1 。

\subsection{题目3}
3．已知 \(V\) 是实数域上次数小于等于 3 的多形式组成的线性空间，当 \(a=\) \(\_\_\_\_\)时，该线性空间中满足 \(f(1)=a, f(-1)=0\) 的多项式集合是 \(V\) 的子空间，关于内积 \((f, g)=\int_{-1}^{1} f(x) g(x) \mathrm{d} x\) 下该子空间的单位正交基为 \(\_\_\_\_\)。

\subsection{题目4}
4．已知 \(f_{1}=2+2 x^{2}+3 x^{3}, f_{2}=1+x+x^{2}+3 \dot{x}^{3}, f_{3}=5+x+5 x^{2}+9 x^{3}, f_{4}=2 x+3 x^{3}, W\) 是由它们张成的线性空间，那么从基 \(f_{1}, f_{2}\) 到基 \(f_{3}, f_{4}\) 的过渡矩阵为 \(\_\_\_\_\)，向量组 \(f_{1}, f_{3}-f_{2}, 2 f_{3}+f_{4}\) 的秩等于 \(\_\_\_\_\)。

\begin{itemize}

\subsection{题目5}
\end{itemize}

\subsection{题目5}
原题边界“5．”：设矩阵 \(A\) 的秩等于 \(3, b\) 不等于零，\(A X=b\) 有解
\[
X_{1}=(1,-1,2,3,1)^{\mathrm{T}}, X_{2}=(0,1,-1,0,-2)^{\mathrm{T}}, X_{3}=(-1,1,2,1,3)^{\mathrm{T}} .
\]
那么 \(A X=b\) 的通解用 \(X_{1}, X_{2}, X_{3}\) 表示为 \(\_\_\_\_\)，给出一个所有解集合的极大线性无关组 \(\_\_\_\_\)．
\end{document}
